\section{回归模型的选择与预测区间}
\label{sec:regression-selection}

前面各模型回答不同问题：均值模型预测条件平均，分位模型预测阈值，贝叶斯与GPR给分布与区间。比较时先对齐交付对象，再走完整的拟合、选择、校准、测试流程，最后才谈覆盖与评分。

\subsection{比较目标与完整预测流程}
\label{sec:regression-comparison-goal}

先确定要交付条件均值、指定分位点、容忍误差最优预测还是响应分布，分别选择匹配的损失与检查对象：条件均值比较测试均方误差，分位数比较相应损失与经验超越率，未来响应区间同时报告覆盖率、宽度与分组表现；不能用MSE评价分位模型，也不能用平均宽度掩盖覆盖率不足。模型家族名称不能替代目标说明，给同一房价任务的均值与预算目标作对照。各类输出在拟合、选择、校准与最终测试中的分工，以及对应指标，见表\ref{tab:regression-evaluation-protocol}。

在训练范围内共同拟合预处理、表示、参数和超参数，区分其中各层优化；使用验证或嵌套验证比较完整过程，不预先在全部数据上选成分、结点或变量。固定相同数据划分、目标和预算，比较泛化误差、解释、存储和新输入成本；局部方法重拟合与树路径、核展开等代价分别计入，训练目标低、目标下降或后验计算精确，各自只能支持相应结论。若还需要预测区间，先区分模型内预测分布与校准覆盖，再决定是否保留独立校准数据；选择和校准完成后冻结过程，在最终测试数据评价既定目标，不能重复利用测试结果选择配置。用来报告最终水平的数据不能在训练或挑选时被看过：拟合与超参数选择用训练折和验证折，共形区间另用独立校准记录，最后才用完全没参与选择的测试集报告。

% 本章共同的输出评价；原连续参数风险模板已按本轮模型讲解取舍移除。
各类回归输出在拟合、选择、校准与最终测试中的分工见表\ref{tab:regression-evaluation-protocol}。

\begin{table*}[htbp]
  \centering
  \small
  \begin{tabularx}{\linewidth}{@{}>{\RaggedRight\arraybackslash}p{22mm}>{\RaggedRight\arraybackslash}X>{\RaggedRight\arraybackslash}X>{\RaggedRight\arraybackslash}X@{}}
    \toprule
    输出 & 拟合与验证选择 & 独立校准 & 最终测试 \\
    \midrule
    条件均值 & 平方损失或均值模型似然；验证MSE & 点预测不另设 & MSE、残差与分组误差 \\
    条件分位数 & 对应$\tau$的损失；检查尾部记录与交叉 & 原始分位点不另设 & 分位数损失与经验超越率 \\
    模型内响应预测区间 & 拟合响应或后验模型；选择超参数 & 按模型构造；若再校准则单独留出 & 响应覆盖率、宽度与分组表现 \\
    共形覆盖区间 & 拟合基础预测器；校准前选定模型 & 留出记录估计不合格分数分位数 & 覆盖率、宽度与分组表现 \\
    完整分布 & 似然或适当评分；验证分布诊断 & 若作再校准则单独留出 & 负对数得分或CRPS、校准与锐度 \\
    \bottomrule
  \end{tabularx}
  \caption{不同回归输出的评价分工。两类区间均以未来响应为覆盖对象：模型内响应预测区间依赖其响应模型与参数估计或后验设定，共形覆盖区间用独立校准记录估计覆盖修正。最终测试留到模型选择与校准结束之后。}
  \label{tab:regression-evaluation-protocol}
\end{table*}

评价应与交付对象对应。在逐个比较指标之前，先要保证一件事：用来报告最终水平的数据，不能在模型训练或挑选时就被看过。因此一条记录在整个流程中要分工明确——拟合与选择超参数用训练折和验证折，共形区间另用独立校准记录，最后才用完全没参与任何选择的测试集报告结果；表\ref{tab:regression-evaluation-protocol}把这套分工按交付对象列清。在此基础上，条件均值比较测试均方误差（MSE），即残差平方的平均；分位数比较相应损失及经验超越率；未来响应区间同时报告覆盖率、宽度和分组表现。参数区间与潜在函数值区间覆盖各自的真值，不能直接拿带噪响应当条件均值真值。GP预测响应时应使用$v_\star+\sigma^2$，而不是只用函数方差$v_\star$。

完整分布可用\term{适当评分}（proper scoring rule）比较：在相关期望存在时，报告真实分布能取得最优期望评分。统一按损失越小越好计，负对数得分为$-\log p(y)$；给实际响应分配的概率或密度越低，损失越大。连续密度可以超过$1$，因此该得分不必非负。连续秩概率得分（continuous ranked probability score, CRPS）比较累积分布$F$与观测所给的阶跃函数：
\begin{equation}
 \operatorname{CRPS}(F,y)=\int_{\mathbb R}
        [F(t)-\mathbf1\{y\leq t\}]^2\,dt.
 \label{eq:regression-crps}
\end{equation}
有限一阶矩保证此积分有限，并保留响应单位；若$F$只在单点$a$取值，两个阶跃函数仅在$a$和$y$之间不同，积分就等于$|a-y|$。同均值、不同分散程度的分布可以获得不同评分，均值MSE不能替代整条分布的评价\scite{gneiting2007scoring}。

在房价任务中，先按部署时间或地区划分数据，每个训练折内完成缺失值处理、标准化、特征选择与拟合，再用验证结果选择表示和超参数。若需要共形覆盖区间，应在模型冻结后用独立校准记录确定修正，最后才用未参与选择或校准的测试集报告结果。看过测试结果后再调模型，该批数据便不再是最终测试证据。模型内区间是否需要另作校准，取决于交付要求，不由“报告区间”自动决定。


\subsection{预测区间与独立校准}
\label{sec:regression-prediction-interval}

此处是对已训练预测器的区间构造，不另造回归模型；共形校准是接收冻结预测器的独立流程，不挂在任何单个模型之下。

\subsubsection{区间形式与校准计算}
\label{sec:regression-conformal}

\paragraph{区间概率的解释}

点预测不给把握。区分模型内后验概率与固定真实机制下的频率覆盖，说明二者不能只按区间宽窄比较。

\paragraph{分割共形校准}

\term{回归共形预测}（conformal prediction）在冻结预测器之外，用独立校准集给出有限样本边际覆盖保证的区间。分割共形（split conformal）把数据分为训练集与校准集：训练集拟合模型，校准集只用于计算非一致性分数。按冻结预测器、留出残差、校准分位点、形成新输入区间的步骤构造：对回归取残差分数$r_i=|y_i-\widehat f(\bm{x}_i)|$（或标准化后的分数），取校准集$N$个分数的水平$\lceil(N+1)(1-\alpha)/N\rceil$经验分位$Q_{1-\alpha}$，新输入区间为
\[
  [\widehat f(\bm{x}_\star)-Q_{1-\alpha},\ \widehat f(\bm{x}_\star)+Q_{1-\alpha}].
\]
解释$N+1$排序与无穷区间的边界；在交换性（校准与新样本同分布、可交换）条件下，该区间以至少$1-\alpha$的边际概率覆盖未来响应\scite{angelopoulos2022conformal}。

\subsubsection{覆盖性质与使用条件}
\label{sec:regression-coverage}

分割共形保证的是边际覆盖：在全体新样本上平均覆盖率至少$1-\alpha$，不保证每个局部区域都达到该覆盖率；条件覆盖需要更强假设或局部校准。分位模型给出的两个端点可由\term{分布校准回归}（CQR）用共形分数同时上移、下移，使上下限各自满足覆盖\scite{romano2019cqr}。区间宽度与覆盖率相互权衡：覆盖越宽越稳妥，却失去决策价值。除覆盖率与平均宽度外，完整分布用CRPS比较，见式\eqref{eq:regression-crps}；评分应与交付对象对应，训练、校准、测试分工不变。完整覆盖理论指向卷三，本章重点是接到已训练回归器上。
